\textbf{Phần 1. Trắc nghiệm lựa chọn}
\setcounterpageref{ex}{0}

\begin{ex} %[5H3N17]
Một viên gạch men hình vuông có cạnh dài $50\text{ cm}$. Diện tích của viên gạch men đó là
\choice
{$200\text{ cm}^2$}
{$2\,500\text{ cm}^2$}
{\True $0{,}25\text{ m}^2$}
{$2{,}5\text{ m}^2$}
\loigiai{
Diện tích viên gạch là: $50 \times 50 = 2\,500\text{ (cm}^2\text{)} = 0{,}25\text{ m}^2$.\\
\textbf{Đáp án của câu hỏi là C}
}
\end{ex}

\begin{ex} %[5H3N17]
Mỗi hộp gạch men chứa $4$ viên gạch hình vuông cạnh $50\text{ cm}$. Diện tích gạch trong mỗi hộp lát được là
\choice
{\True $1\text{ m}^2$}
{$0{,}5\text{ m}^2$}
{$2\text{ m}^2$}
{$4\text{ m}^2$}
\loigiai{
Diện tích một viên gạch là $0{,}25\text{ m}^2$. Bốn viên gạch có diện tích là: $4 \times 0{,}25 = 1\text{ (m}^2\text{)}$.\\
\textbf{Đáp án của câu hỏi là A}
}
\end{ex}

\begin{ex} %[5H3H17]
Một tấm gỗ hình chữ nhật có kích thước $200\text{ cm} \times 25\text{ cm}$. Diện tích tấm gỗ đó là
\choice
{$5\,000\text{ m}^2$}
{$50\text{ m}^2$}
{\True $0{,}5\text{ m}^2$}
{$5\text{ m}^2$}
\loigiai{
Đổi $200\text{ cm} = 2\text{ m}$; $25\text{ cm} = 0{,}25\text{ m}$.\\
Diện tích tấm gỗ là: $2 \times 0{,}25 = 0{,}5\text{ (m}^2\text{)}$.\\
\textbf{Đáp án của câu hỏi là C}
}
\end{ex}

\begin{ex} %[5H3H17]
Một tấm gỗ nhỏ có kích thước $40\text{ cm} \times 25\text{ cm}$. Diện tích tấm gỗ đó theo đơn vị mét vuông là
\choice
{\True $0{,}1\text{ m}^2$}
{$1\text{ m}^2$}
{$0{,}01\text{ m}^2$}
{$10\text{ m}^2$}
\loigiai{
Đổi $40\text{ cm} = 0{,}4\text{ m}$; $25\text{ cm} = 0{,}25\text{ m}$.\\
Diện tích tấm gỗ là: $0{,}4 \times 0{,}25 = 0{,}1\text{ (m}^2\text{)}$.\\
\textbf{Đáp án của câu hỏi là A}
}
\end{ex}

\begin{ex} %[5H3H17]
Các bạn lớp 5A cần mua $4$ tấm gỗ lớn (mỗi tấm diện tích $0{,}5\text{ m}^2$) và $10$ tấm gỗ nhỏ (mỗi tấm diện tích $0{,}1\text{ m}^2$) để đóng kệ sách. Tổng diện tích gỗ cần mua là
\choice
{$2\text{ m}^2$}
{$4\text{ m}^2$}
{$5\text{ m}^2$}
{\True $3\text{ m}^2$}
\loigiai{
Tổng diện tích gỗ là: $4 \times 0{,}5 + 10 \times 0{,}1 = 2 + 1 = 3\text{ (m}^2\text{)}$.\\
\textbf{Đáp án của câu hỏi là D}
}
\end{ex}

\begin{ex} %[5H3H17]
Giá tiền $1\text{ m}^2$ gỗ công nghiệp là $250\,000$ đồng. Để mua $3\text{ m}^2$ gỗ đó làm giá sách, các bạn cần chi trả số tiền là
\choice
{$500\,000$ đồng}
{$1\,000\,000$ đồng}
{\True $750\,000$ đồng}
{$1\,250\,000$ đồng}
\loigiai{
Số tiền mua gỗ là: $3 \times 250\,000 = 750\,000$ (đồng).\\
\textbf{Đáp án của câu hỏi là C}
}
\end{ex}

\begin{ex} %[5H3H17]
Một phòng học hình chữ nhật có chiều dài $8\text{ m}$ và chiều rộng $6\text{ m}$. Diện tích sàn phòng học đó là
\choice
{$28\text{ m}^2$}
{\True $48\text{ m}^2$}
{$40\text{ m}^2$}
{$56\text{ m}^2$}
\loigiai{
Diện tích mặt sàn phòng học là: $8 \times 6 = 48\text{ (m}^2\text{)}$.\\
\textbf{Đáp án của câu hỏi là B}
}
\end{ex}

\begin{ex} %[5H3V17]
Để lát kín mặt sàn phòng học diện tích $48\text{ m}^2$ bằng các hộp gạch (mỗi hộp có $4$ viên gạch hình vuông cạnh $50\text{ cm}$, lát được $1\text{ m}^2$), người ta cần mua số hộp gạch là
\choice
{\True $48$ hộp}
{$12$ hộp}
{$24$ hộp}
{$96$ hộp}
\loigiai{
Vì mỗi hộp gạch lát được $1\text{ m}^2$ sàn nên cần: $48 : 1 = 48$ (hộp gạch).\\
\textbf{Đáp án của câu hỏi là A}
}
\end{ex}

\begin{ex} %[5H3V17]
Mỗi hộp gạch men có giá $140\,000$ đồng. Chi phí mua $48$ hộp gạch để lát một phòng học là
\choice
{$5\,600\,000$ đồng}
{$7\,200\,000$ đồng}
{$8\,400\,000$ đồng}
{\True $6\,720\,000$ đồng}
\loigiai{
Chi phí mua gạch là: $48 \times 140\,000 = 6\,720\,000$ (đồng).\\
\textbf{Đáp án của câu hỏi là D}
}
\end{ex}

\begin{ex} %[5H3V17]
Trường Tiểu học Đoàn Kết có $10$ phòng học cùng kích thước $8\text{ m} \times 6\text{ m}$. Tổng chi phí mua gạch để lát lại toàn bộ sàn $10$ phòng học đó (biết chi phí lát mỗi phòng là $6\,720\,000$ đồng) là
\choice
{$60\,720\,000$ đồng}
{\True $67\,200\,000$ đồng}
{$672\,000\,000$ đồng}
{$58\,400\,000$ đồng}
\loigiai{
Tổng chi phí lát $10$ phòng học là: $10 \times 6\,720\,000 = 67\,200\,000$ (đồng).\\
\textbf{Đáp án của câu hỏi là B}
}
\end{ex}

\begin{ex} %[5H3V17]
Một mảnh sân hình chữ nhật dài $15\text{ m}$ và rộng $8\text{ m}$. Người ta lát sân bằng loại gạch hình vuông có diện tích $0{,}25\text{ m}^2$/viên. Số viên gạch cần dùng là
\choice
{$120$ viên}
{$240$ viên}
{$960$ viên}
{\True $480$ viên}
\loigiai{
Diện tích mảnh sân là: $15 \times 8 = 120\text{ (m}^2\text{)}$.\\
Số viên gạch cần dùng là: $120 : 0{,}25 = 480$ (viên).\\
\textbf{Đáp án của câu hỏi là D}
}
\end{ex}

\begin{ex} %[5H3V17]
Một bức tường hình chữ nhật có diện tích $18\text{ m}^2$. Người ta ốp tường bằng các tấm gỗ có kích thước $60\text{ cm} \times 10\text{ cm}$. Số tấm gỗ cần dùng để ốp kín bức tường là
\choice
{$30$ tấm}
{\True $300$ tấm}
{$60$ tấm}
{$150$ tấm}
\loigiai{
Đổi diện tích mỗi tấm gỗ: $0{,}6\text{ m} \times 0{,}1\text{ m} = 0{,}06\text{ m}^2$.\\
Số tấm gỗ cần dùng là: $18 : 0{,}06 = 300$ (tấm).\\
\textbf{Đáp án của câu hỏi là B}
}
\end{ex}

\textbf{Phần 2. Trắc nghiệm đúng sai}
\setcounterpageref{ex}{0}

\begin{ex} %[5H3H17]
Các bạn học sinh lớp 5A dự định đóng kệ sách gồm $4$ tấm gỗ lớn kích thước $200\text{ cm} \times 25\text{ cm}$ và $10$ tấm gỗ nhỏ kích thước $40\text{ cm} \times 25\text{ cm}$. Giá gỗ là $250\,000$ đồng/$1\text{ m}^2$. Xét tính đúng sai của các khẳng định sau:
\choiceTF
{\True Diện tích của $4$ tấm gỗ lớn là $2\text{ m}^2$.}
{Diện tích của $10$ tấm gỗ nhỏ là $0{,}5\text{ m}^2$.}
{\True Tổng diện tích tất cả các tấm gỗ cần dùng là $3\text{ m}^2$.}
{\True Số tiền cần để mua gỗ đóng chiếc kệ sách là $750\,000$ đồng.}
\loigiai{
\begin{itemize}
    \item Diện tích $1$ tấm gỗ lớn: $2 \times 0{,}25 = 0{,}5\text{ m}^2 \Rightarrow 4$ tấm lớn có diện tích $4 \times 0{,}5 = 2\text{ m}^2$.
    \item Diện tích $1$ tấm gỗ nhỏ: $0{,}4 \times 0{,}25 = 0{,}1\text{ m}^2 \Rightarrow 10$ tấm nhỏ có diện tích $10 \times 0{,}1 = 1\text{ m}^2 \neq 0{,}5\text{ m}^2$.
    \item Tổng diện tích gỗ: $2 + 1 = 3\text{ m}^2$.
    \item Chi phí mua gỗ: $3 \times 250\,000 = 750\,000$ đồng.
\end{itemize}
\begin{itemize}
    \item a) \textbf{Đúng}.
    \item b) \textbf{Sai}: Diện tích là $1\text{ m}^2$.
    \item c) \textbf{Đúng}.
    \item d) \textbf{Đúng}.
\end{itemize}
}
\end{ex}

\begin{ex} %[5H3V17]
Một phòng học có chiều dài $8\text{ m}$, chiều rộng $6\text{ m}$ được lát sàn bằng gạch vuông cạnh $50\text{ cm}$ (đóng hộp $4$ viên/hộp với giá $140\,000$ đồng/hộp). Xét tính đúng sai của các khẳng định sau:
\choiceTF
{\True Diện tích sàn phòng học là $48\text{ m}^2$.}
{Diện tích gạch trong mỗi hộp lát được $0{,}5\text{ m}^2$ sàn.}
{\True Cần mua tất cả $48$ hộp gạch để lát kín phòng học.}
{\True Chi phí mua gạch để lát một phòng học là $6\,720\,000$ đồng.}
\loigiai{
\begin{itemize}
    \item Diện tích phòng học: $8 \times 6 = 48\text{ m}^2$.
    \item Mỗi viên gạch có diện tích $0{,}5 \times 0{,}5 = 0{,}25\text{ m}^2 \Rightarrow 1$ hộp có $4 \times 0{,}25 = 1\text{ m}^2 \neq 0{,}5\text{ m}^2$.
    \item Số hộp gạch cần mua: $48 : 1 = 48$ hộp.
    \item Chi phí mua gạch: $48 \times 140\,000 = 6\,720\,000$ đồng.
\end{itemize}
\begin{itemize}
    \item a) \textbf{Đúng}.
    \item b) \textbf{Sai}: Mỗi hộp lát được $1\text{ m}^2$.
    \item c) \textbf{Đúng}.
    \item d) \textbf{Đúng}.
\end{itemize}
}
\end{ex}

\textbf{Phần 3. Trả lời ngắn}
\setcounterpageref{ex}{0}

\begin{ex} %[5H3H17]
Tính diện tích của một tấm gỗ hình chữ nhật có chiều dài $200\text{ cm}$ và chiều rộng $25\text{ cm}$ theo đơn vị mét vuông. (Nhập kết quả dưới dạng số thập phân).
\shortans{0,5}
\loigiai{
\textbf{Đáp số:} $0{,}5$\\
$2\text{ m} \times 0{,}25\text{ m} = 0{,}5\text{ m}^2$.
}
\end{ex}

\begin{ex} %[5H3H17]
Tính tổng diện tích (theo mét vuông) của $10$ tấm gỗ hình chữ nhật cùng kích thước $40\text{ cm} \times 25\text{ cm}$.
\shortans{1}
\loigiai{
\textbf{Đáp số:} $1$\\
Diện tích $10$ tấm gỗ là: $10 \times (0{,}4 \times 0{,}25) = 1\text{ (m}^2\text{)}$.
}
\end{ex}

\begin{ex} %[5H3H17]
Để mua $3\text{ m}^2$ gỗ công nghiệp với đơn giá $250\,000$ đồng/$1\text{ m}^2$ thì phải trả bao nhiêu nghìn đồng?
\shortans{750}
\loigiai{
\textbf{Đáp số:} $750$\\
Số tiền cần trả là: $3 \times 250\,000 = 750\,000$ đồng $= 750$ nghìn đồng.
}
\end{ex}

\begin{ex} %[5H3V17]
Phòng học hình chữ nhật dài $8\text{ m}$, rộng $6\text{ m}$. Cần mua bao nhiêu hộp gạch (mỗi hộp gồm $4$ viên gạch vuông cạnh $50\text{ cm}$) để lát kín nền phòng học?
\shortans{48}
\loigiai{
\textbf{Đáp số:} $48$\\
Diện tích phòng: $8 \times 6 = 48\text{ m}^2$. Mỗi hộp gạch lát được $4 \times 0{,}25 = 1\text{ m}^2 \Rightarrow$ Cần $48$ hộp.
}
\end{ex}

\begin{ex} %[5H3V17]
Mỗi hộp gạch có giá $140\,000$ đồng. Hỏi cần bao nhiêu triệu đồng để mua gạch lát nền cho $5$ phòng học, biết mỗi phòng cần $48$ hộp gạch? (Nhập kết quả dưới dạng số thập phân).
\shortans{33,6}
\loigiai{
\textbf{Đáp số:} $33{,}6$\\
Chi phí lát $5$ phòng: $5 \times 48 \times 140\,000 = 33\,600\,000$ đồng $= 33{,}6$ triệu đồng.
}
\end{ex}

\begin{ex} %[5H3V17]
Một bức tường hình chữ nhật diện tích $18\text{ m}^2$ được ốp kín bằng các tấm gỗ nhỏ diện tích $0{,}06\text{ m}^2$/tấm. Hỏi cần tất cả bao nhiêu tấm gỗ nhỏ như vậy?
\shortans{300}
\loigiai{
\textbf{Đáp số:} $300$\\
Số tấm gỗ cần dùng là: $18 : 0{,}06 = 300$ (tấm).
}
\end{ex}
